\section{Discussion}\label{sec:discussion}

The book is not closed on the design of straight-line asynchrony. Communities are actively exploring extensions to the core async/await semantics presented in this work.

\citet{elizarov_kotlin_2021} use the \emph{suspend} keyword in Kotlin to define async functions. Instead of an explicit await keyword, the language automatically introduces suspension points where necessary. They argue that this approach avoids forgotten awaits, and further that it improves the look-and-feel of calling a suspending function. However, this claim, and other consequences of this design, are not accompanied by an evaluation.

The Zig programming language is stabilizing its ``colorblind'' async/await~\cite{kelley_zigs_2025,cro_what_2020}. Functions are parameterized over an \abbr{i/o} implementation; libraries exploit asynchrony when available, and clients pass whichever \abbr{i/o} implementation they have, including a synchronous variant. This approach facilitates sync--async interop, but introduces new risks where libraries written explicitly for asynchrony deadlock or panic at runtime if asynchrony is unavailable.

\Cpp{}20 provides a flexible implementation of coroutines that does not ascribe a semantics to its async/await~\cite{isoiec_jtc1_sc22_wg21_technical_2017}. We cannot attribute \Cpp to any particular design point in the taxonomy provided in \Cref{tab:dims-overview} because each axis is configurable. Although elegant and neutral, the choice of full programmability makes each library an async \abbr{dsl}; knowledge transfer between projects within the same language becomes exceedingly difficult.

\paragraph{Where do we go from here?}

We set out to examine the space of straight-line asynchrony, ourselves confused by some statements we read in individual language descriptions. In particular, as we note in \Cref{s:sync-quotes-from-langs}, these languages motivate async from a perspective of similarity to synchrony, which we found difficult to reconcile. Our design space exploration demonstrates the numerous ways in which this analogy fails.

The principal use for our exploration is to help developers. Developers transitioning from one language to another carry their prior knowledge to make sense of the new in terms of the old. It is especially confusing when the new language has the same keywords and similar description, but different semantics. We believe that our articulation of the design space, such as \Cref{tab:dims-overview}, can help learners and educators talk more effectively about how to translate knowledge between languages.
%
We also believe that our exploration can help language designers build new features for straight-line asynchrony. Our design space can help enumerate key design dimensions, preventing important details from slipping through the cracks. It can also help designers more effectively compare and contrast their design against related languages.

An open challenge is to build a body of empirical data that can support designers in reasoning about the trade-offs of different approaches. For example, how much overhead does \seager async function application incur versus \eager, and how much more concurrency does it gain? To what extent does \simultaneous cancellation prevent realistic bugs that would occur with a \bup cancellation strategy? Which of these semantics best matches developer expectations for the behavior of asynchronous code? Which decisions lead to more errors? We hope that this exploration can inspire practitioners and academics alike to build a better understanding of the evolving world of asynchronous programming.
